\documentclass[aspectratio=169,10pt]{beamer}

% Shared Beamer setup for Fall 2026 course slides.
% Clean Cal Poly-inspired theme: no custom class files or uncommon packages.
\mode<presentation>{
  \usetheme{default}
  \usefonttheme{professionalfonts}
}

\definecolor{CalPolyGreen}{HTML}{154734}
\definecolor{CalPolyGold}{HTML}{BD8B13}
\definecolor{CalPolySage}{HTML}{7FA28F}
\definecolor{CalPolyMint}{HTML}{DDF1E7}
\definecolor{CalPolySand}{HTML}{A29F86}
\definecolor{CalPolyBeige}{HTML}{EFEEDE}
\definecolor{DeckInk}{HTML}{1F2933}
\definecolor{DeckMuted}{HTML}{5F6B7A}
\definecolor{DeckSoft}{HTML}{F7F7F5}

\setbeamersize{text margin left=0.55in,text margin right=0.55in}
\setbeamercolor{normal text}{fg=DeckInk,bg=white}
\setbeamercolor{structure}{fg=CalPolyGreen}
\setbeamercolor{title}{fg=white}
\setbeamercolor{frametitle}{fg=CalPolyGreen,bg=white}
\setbeamercolor{block title}{bg=CalPolySage,fg=white}
\setbeamercolor{block body}{bg=CalPolyMint,fg=black}
\setbeamercolor{alerted text}{fg=CalPolyGold}

\setbeamerfont{title}{size=\Huge,series=\bfseries}
\setbeamerfont{frametitle}{size=\Large,series=\bfseries}
\setbeamerfont{block title}{size=\normalsize,series=\bfseries}

\setbeamertemplate{navigation symbols}{}
\setbeamertemplate{itemize item}{\color{CalPolyGold}$\blacktriangleright$}
\setbeamertemplate{itemize subitem}{\color{CalPolyGreen}$\bullet$}
\setbeamertemplate{footline}{
  \hfill{\color{DeckMuted}\scriptsize\insertframenumber/\inserttotalframenumber}
  \hspace{0.18in}\vspace{0.12in}
}
\setbeamertemplate{frametitle}{
  \vspace{0.18in}
  \hspace{0.02in}\insertframetitle\par
  \vspace{0.08in}
  \noindent\makebox[\linewidth][l]{%
    {\color{CalPolyGold}\rule{0.55in}{2pt}}%
    {\color{CalPolyGreen}\rule{\dimexpr\linewidth-0.55in\relax}{2pt}}%
  }
}

\newcommand{\CourseNumber}{Course}
\newcommand{\CourseTitle}{Course Title}
\newcommand{\LectureNumber}{0}
\newcommand{\LectureTitle}{Lecture Title}
\newcommand{\LectureDate}{Date}
\newcommand{\CourseUnit}{Unit}

\newcommand{\coursenumber}[1]{\renewcommand{\CourseNumber}{#1}}
\newcommand{\coursetitle}[1]{\renewcommand{\CourseTitle}{#1}}
\newcommand{\lecturenumber}[1]{\renewcommand{\LectureNumber}{#1}}
\newcommand{\lecturetitle}[1]{\renewcommand{\LectureTitle}{#1}}
\newcommand{\lecturedate}[1]{\renewcommand{\LectureDate}{#1}}
\newcommand{\courseunit}[1]{\renewcommand{\CourseUnit}{#1}}

\author{Kirk Duran}
\institute{California Polytechnic State University, San Luis Obispo}
\date{\LectureDate}

\newcommand{\makedecktitle}{
  \begingroup
  \setbeamercolor{background canvas}{bg=CalPolyGreen}
  \begin{frame}[plain]
    \vfill
    \begin{center}
      {\usebeamerfont{title}\color{white}\LectureTitle\par}
      \vspace{0.18in}
      {\Large\color{white}Lecture \LectureNumber\par}
    \end{center}
    \vfill
    \noindent{\color{white}\rule{\linewidth}{1.2pt}}\par
    \vspace{0.08in}
    \noindent\colorbox{CalPolyGold}{\makebox[0.24\linewidth][c]{\color{white}\Large\CourseNumber}}
    \hspace{0.12in}{\color{white}\Large Kirk Duran}
  \end{frame}
  \endgroup
}

\newenvironment{keyidea}{\begin{block}{Key idea}}{\end{block}}
\newenvironment{activity}{\begin{block}{In class}}{\end{block}}
\newenvironment{cpdefinition}{\begin{block}{Definition}}{\end{block}}
\newenvironment{exampleblockcp}{
  \begingroup
  \setbeamercolor{block title}{bg=CalPolySand,fg=white}
  \setbeamercolor{block body}{bg=CalPolyBeige,fg=black}
  \begin{block}{Example}
}{
  \end{block}
  \endgroup
}

\newcommand{\imageplaceholder}[2][0.3in]{%
  \noindent\makebox[\linewidth][c]{%
    \fcolorbox{CalPolySage}{CalPolyBeige}{%
      \parbox[c][#1][c]{0.86\linewidth}{%
        \centering
        {\color{DeckMuted}\scriptsize Image placeholder: }%
        {\color{DeckInk}\scriptsize\ttfamily\detokenize{#2}}%
      }%
    }%
  }%
}

\newcommand{\sectionframe}[1]{
  \begingroup
  \setbeamercolor{background canvas}{bg=CalPolyGreen}
  \begin{frame}[plain]
    \vfill
    {\color{white}\LARGE\bfseries #1\par}
    \vspace{0.18in}
    {\color{CalPolyGold}\rule{0.22\paperwidth}{3pt}}
    {\color{white}\rule{0.62\paperwidth}{3pt}}
    \vfill
  \end{frame}
  \endgroup
}

\newcommand{\placeholdernote}{
  \begin{block}{Placeholder}
    Replace these bullets with the final lecture material, examples, and in-class activities.
  \end{block}
}

\usepackage[T1]{fontenc}
\usepackage{lmodern}
\pdfstringdefDisableCommands{\def\translate#1{#1}}

\graphicspath{{../assets/lecture-28/}}

% Generated classroom visuals live in ../assets/lecture-28/.
\newcommand{\lectureimage}[2][1.75in]{%
  \includegraphics[width=\linewidth,height=#1,keepaspectratio]{#2}%
}

\setbeamersize{text margin left=0.42in,text margin right=0.42in}

\coursenumber{CS 1001}
\coursetitle{Introduction to Computing}
\lecturenumber{28}
\lecturetitle{Objects I --- Compound Data}
\lecturedate{Fall 2026}
\courseunit{11. Objects and Compound Data}

\begin{document}

\makedecktitle

\begin{frame}{Today}
  \begin{itemize}
    \item Explain why related values sometimes need to represent one conceptual entity.
    \item Define \textbf{class}, \textbf{instance}, \textbf{attribute}, and \textbf{object state}.
    \item Define a simple class and initialize instance attributes with \texttt{\_\_init\_\_}.
    \item Construct multiple independent instances from the same class definition.
    \item Read and update attributes using dot notation.
    \item Extend our existing name--reference model from lists to user-defined objects.
    \item Compare an object representation with a dictionary representation of the same data.
  \end{itemize}

  \begin{keyidea}
    An object lets several related values collectively represent one thing; a class describes the structure shared by objects of that kind.
  \end{keyidea}
\end{frame}

\begin{frame}{The Boundary for This Lecture}
  \begin{columns}[T,onlytextwidth]
    \begin{column}{0.47\textwidth}
      \begin{block}{Today}
        \begin{itemize}
          \item compound data
          \item classes and instances
          \item construction and initialization
          \item instance attributes
          \item object state
          \item references to objects
          \item object vs. dictionary representation
        \end{itemize}
      \end{block}
    \end{column}
    \begin{column}{0.47\textwidth}
      \begin{alertblock}{Next lecture}
        \begin{itemize}
          \item general methods
          \item deeper treatment of \texttt{self}
          \item mutation through methods
          \item object aliasing
          \item identity and equality
          \item lists of objects
          \item objects containing collections
        \end{itemize}
      \end{alertblock}
    \end{column}
  \end{columns}
\end{frame}

\sectionframe{Part 1: Why Compound Data?}

\begin{frame}{One Entity Can Require Several Related Values}
  Suppose a program needs to represent two cars.

  \begin{block}{Separate variables}
    \texttt{car1\_make = "Honda"}\\
    \texttt{car1\_model = "Civic"}\\
    \texttt{car1\_year = 2020}\\[0.08in]
    \texttt{car2\_make = "Toyota"}\\
    \texttt{car2\_model = "Corolla"}\\
    \texttt{car2\_year = 2022}
  \end{block}

  \begin{itemize}
    \item The program contains the necessary values.
    \item Their relationship exists only in the variable names.
    \item Adding another car requires another group of coordinated names.
  \end{itemize}
\end{frame}

\begin{frame}{A List Groups the Values but Hides Their Roles}
  \begin{columns}[T,onlytextwidth]
    \begin{column}{0.54\textwidth}
      \begin{block}{List representation}
        \texttt{car1 = ["Honda", "Civic", 2020]}
      \end{block}

      \begin{itemize}
        \item One value now contains the three related components.
        \item But their meaning depends on remembering a positional convention.
        \item \texttt{car1[0]} means make only because the programmer decided that index 0 means make.
      \end{itemize}
    \end{column}
    \begin{column}{0.40\textwidth}
      % Intended visual: a three-slot list labeled 0,1,2 contrasted with make/model/year labels.
      \lectureimage[2.22in]{car-list-positional-fields.png}
    \end{column}
  \end{columns}
\end{frame}

\begin{frame}{A Dictionary Gives the Components Meaningful Names}
  \begin{block}{Dictionary representation}
    \texttt{car1 = \{"make": "Honda", "model": "Civic", "year": 2020\}}
  \end{block}

  \begin{columns}[T,onlytextwidth]
    \begin{column}{0.47\textwidth}
      \begin{block}{Access}
        \texttt{car1["make"]}\\
        \texttt{car1["model"]}\\
        \texttt{car1["year"]}
      \end{block}
    \end{column}
    \begin{column}{0.47\textwidth}
      \begin{itemize}
        \item The keys document each component's role.
        \item The values are still grouped into one structure.
        \item This is often an excellent representation for associative data.
      \end{itemize}
    \end{column}
  \end{columns}
\end{frame}

\begin{frame}{Definition: Compound Data}
  \begin{cpdefinition}
    \textbf{Compound data} represents one conceptual value using multiple related component values.
  \end{cpdefinition}

  \begin{itemize}
    \item A date might contain a year, month, and day.
    \item A student might contain a name, major, and GPA.
    \item A point might contain an $x$ coordinate and a $y$ coordinate.
    \item A car might contain a make, model, and year.
  \end{itemize}

  \begin{block}{New question}
    How can we define a reusable structure that says what components every \emph{Car} should contain?
  \end{block}
\end{frame}

\begin{frame}{Objects Combine Related State Under One Identity}
  \begin{columns}[T,onlytextwidth]
    \begin{column}{0.54\textwidth}
      \begin{cpdefinition}
        An \textbf{object} is a value with its own identity and state. Its state can contain multiple related attributes.
      \end{cpdefinition}

      \begin{itemize}
        \item \textbf{Attributes} store information associated with the object.
        \item Objects can also have behavior through methods.
        \item Today we focus on the \textbf{state} side of this model.
      \end{itemize}
    \end{column}
    \begin{column}{0.40\textwidth}
      % Intended visual: one Car entity containing labeled make/model/year fields, with behavior faded for next lecture.
      \lectureimage[2.22in]{object-state-and-behavior-overview.png}
    \end{column}
  \end{columns}
\end{frame}

\sectionframe{Part 2: Classes and Instances}

\begin{frame}{A Class Describes a Kind of Object}
  \begin{cpdefinition}
    A \textbf{class} is a program definition that describes the structure and behavior shared by objects of a particular kind.
  \end{cpdefinition}

  \begin{itemize}
    \item The class is a reusable definition.
    \item A class does not describe one particular Honda or Toyota.
    \item It describes what information each \texttt{Car} instance should maintain.
  \end{itemize}

  \begin{keyidea}
    Think \textbf{class = definition of a kind}; \textbf{instance = one particular value of that kind}.
  \end{keyidea}
\end{frame}

\begin{frame}{First Class Definition}
  \begin{block}{Car class}
    \texttt{class Car:}\\
    \hspace*{0.20in}\texttt{def \_\_init\_\_(self, make: str, model: str, year: int):}\\
    \hspace*{0.40in}\texttt{self.make = make}\\
    \hspace*{0.40in}\texttt{self.model = model}\\
    \hspace*{0.40in}\texttt{self.year = year}
  \end{block}

  \begin{itemize}
    \item \texttt{class Car:} begins the definition of a new class.
    \item The indented code belongs to that class definition.
    \item \texttt{\_\_init\_\_} establishes the initial state of a new instance.
    \item Type hints document intended types; Python does not enforce them automatically.
  \end{itemize}
\end{frame}

\begin{frame}{An Instance Is One Particular Object}
  \begin{cpdefinition}
    An \textbf{instance} is an individual object created according to a class definition.
  \end{cpdefinition}

  \begin{block}{Construct two instances}
    \texttt{car1 = Car("Honda", "Civic", 2020)}\\
    \texttt{car2 = Car("Toyota", "Corolla", 2022)}
  \end{block}

  \begin{itemize}
    \item Both objects are instances of \texttt{Car}.
    \item Each instance has its own attribute state.
    \item The variable names \texttt{car1} and \texttt{car2} refer to different instances.
  \end{itemize}
\end{frame}

\begin{frame}{One Class, Many Possible Instances}
  \begin{columns}[T,onlytextwidth]
    \begin{column}{0.50\textwidth}
      \begin{block}{Construction}
        \texttt{car1 = Car("Honda", "Civic", 2020)}\\[0.08in]
        \texttt{car2 = Car("Toyota", "Corolla", 2022)}\\[0.08in]
        \texttt{car3 = Car("Ford", "Mustang", 1967)}
      \end{block}
    \end{column}
    \begin{column}{0.44\textwidth}
      % Intended visual: one Car class definition branching to three distinct instance cards.
      \lectureimage[2.22in]{one-class-many-instances.png}
    \end{column}
  \end{columns}
\end{frame}

\sectionframe{Part 3: Construction and Initialization}

\begin{frame}{Calling the Class Constructs an Instance}
  \begin{block}{Expression}
    \centering
    \Large\texttt{Car("Honda", "Civic", 2020)}
  \end{block}

  This expression produces a new \texttt{Car} instance.

  \begin{enumerate}
    \item Python creates a new instance of the class.
    \item Python invokes the class's \texttt{\_\_init\_\_} method for that instance.
    \item The supplied arguments initialize the instance's attributes.
    \item The resulting instance becomes the value of the expression.
  \end{enumerate}

  \begin{keyidea}
    We construct an object by calling the class. \texttt{\_\_init\_\_} initializes the state of the newly created instance.
  \end{keyidea}
\end{frame}

\begin{frame}{What Does \texttt{self} Mean Here?}
  \begin{columns}[T,onlytextwidth]
    \begin{column}{0.54\textwidth}
      \begin{itemize}
        \item \texttt{self} is a parameter of the method.
        \item During initialization, it refers to the new instance being initialized.
        \item \texttt{self.make} means ``the \texttt{make} attribute belonging to this instance.''
        \item Different constructor calls use different instances as \texttt{self}.
      \end{itemize}

      \begin{alertblock}{For today}
        We only need enough \texttt{self} to understand initialization. Lecture 29 develops the method-call model in detail.
      \end{alertblock}
    \end{column}
    \begin{column}{0.40\textwidth}
      % Intended visual: __init__ frame with self arrow pointing to one new Car instance.
      \lectureimage[2.22in]{self-during-initialization.png}
    \end{column}
  \end{columns}
\end{frame}

\begin{frame}{Attribute Assignment Creates the Instance State}
  Consider one line from \texttt{\_\_init\_\_}:

  \begin{block}{Statement}
    \centering
    \Large\texttt{self.make = make}
  \end{block}

  \begin{enumerate}
    \item Evaluate the right-hand side name \texttt{make}.
    \item Follow \texttt{self} to the current instance.
    \item Associate the attribute name \texttt{make} with that value on the instance.
  \end{enumerate}

  \begin{keyidea}
    The assignment model has not disappeared: evaluate the right-hand side, then update program state. The destination is now an attribute on an object.
  \end{keyidea}
\end{frame}

\begin{frame}{Trace the First Construction}
  \begin{block}{Statement}
    \texttt{car1 = Car("Honda", "Civic", 2020)}
  \end{block}

  \begin{center}
    \renewcommand{\arraystretch}{1.25}
    \begin{tabular}{l|l}
      \textbf{Name during initialization} & \textbf{Value} \\\hline
      \texttt{make} & \texttt{"Honda"} \\
      \texttt{model} & \texttt{"Civic"} \\
      \texttt{year} & \texttt{2020} \\
      \texttt{self} & new \texttt{Car} instance
    \end{tabular}
  \end{center}

  After the assignments, the instance contains three attributes: \texttt{make}, \texttt{model}, and \texttt{year}.
\end{frame}

\sectionframe{Part 4: Attributes and Object State}

\begin{frame}{Dot Notation Accesses an Attribute}
  \begin{block}{Given}
    \texttt{car1 = Car("Honda", "Civic", 2020)}
  \end{block}

  \begin{columns}[T,onlytextwidth]
    \begin{column}{0.47\textwidth}
      \begin{block}{Expressions}
        \texttt{car1.make} $\rightarrow$ \texttt{"Honda"}\\
        \texttt{car1.model} $\rightarrow$ \texttt{"Civic"}\\
        \texttt{car1.year} $\rightarrow$ \texttt{2020}
      \end{block}
    \end{column}
    \begin{column}{0.47\textwidth}
      \begin{itemize}
        \item The expression left of the dot identifies an object.
        \item The name right of the dot identifies one attribute on that object.
      \end{itemize}
    \end{column}
  \end{columns}
\end{frame}

\begin{frame}{Object State Is the Collection of Current Attribute Values}
  \begin{cpdefinition}
    The \textbf{state} of an object is the collection of information currently stored in its attributes.
  \end{cpdefinition}

  \begin{block}{State of \texttt{car1}}
    \texttt{make = "Honda"}, \texttt{model = "Civic"}, \texttt{year = 2020}
  \end{block}

  \begin{itemize}
    \item State can be inspected through attribute access.
    \item For mutable objects, state can change over time.
    \item The object's identity can remain the same while its state changes.
  \end{itemize}
\end{frame}

\begin{frame}{Attribute Assignment Changes State}
  \begin{block}{Before}
    \texttt{car1.year} $\rightarrow$ \texttt{2020}
  \end{block}

  \begin{block}{Statement}
    \texttt{car1.year = 2021}
  \end{block}

  \begin{block}{After}
    \texttt{car1.year} $\rightarrow$ \texttt{2021}
  \end{block}

  \begin{keyidea}
    The name \texttt{car1} still refers to the same instance; one component of that instance's state has changed.
  \end{keyidea}
\end{frame}

\begin{frame}{Mutation and Rebinding Are Still Different}
  \begin{columns}[T,onlytextwidth]
    \begin{column}{0.47\textwidth}
      \begin{block}{Attribute mutation}
        \texttt{car1.year = 2021}
      \end{block}
      The existing instance changes state.
    \end{column}
    \begin{column}{0.47\textwidth}
      \begin{block}{Name rebinding}
        \texttt{car1 = Car("Ford", "Focus", 2018)}
      \end{block}
      \texttt{car1} is bound to a different instance.
    \end{column}
  \end{columns}

  \begin{keyidea}
    This is the same distinction we used for lists: changing a value through a reference is different from changing which value a name refers to.
  \end{keyidea}
\end{frame}

\begin{frame}{Prediction: Which Object Changes?}
  \begin{activity}
    Predict every printed value before tracing the code.
  \end{activity}

  \begin{block}{Code}
    \texttt{first = Car("Honda", "Civic", 2020)}\\
    \texttt{second = Car("Honda", "Civic", 2020)}\\
    \texttt{first.year = 2025}\\[0.06in]
    \texttt{print(first.year)}\\
    \texttt{print(second.year)}
  \end{block}

  \begin{itemize}
    \item How many \texttt{Car} instances were constructed?
    \item Does having the same attribute values make them one instance?
  \end{itemize}
\end{frame}

\begin{frame}{Reveal: Same State Does Not Mean Same Instance}
  The output is:

  \begin{block}{Output}
    \texttt{2025}\\
    \texttt{2020}
  \end{block}

  Two calls to \texttt{Car(...)} constructed two distinct objects, even though their initial attribute values were equal.

  \begin{center}
    \texttt{first} $\rightarrow$ instance A \qquad\qquad
    \texttt{second} $\rightarrow$ instance B
  \end{center}
\end{frame}

\sectionframe{Part 5: References and Identity}

\begin{frame}{Our Earlier Reference Model Now Generalizes}
  Earlier in the course we used a focused vocabulary for mutable lists:

  \begin{center}
    name $\longrightarrow$ reference $\longrightarrow$ list
  \end{center}

  We can now generalize that model:

  \begin{center}
    \Large name $\longrightarrow$ reference $\longrightarrow$ object
  \end{center}

  \begin{keyidea}
    Lists were not a special exception. They are one kind of object. User-defined class instances participate in the same reference model.
  \end{keyidea}
\end{frame}

\begin{frame}{A Variable Stores a Reference to the Instance}
  \begin{columns}[T,onlytextwidth]
    \begin{column}{0.51\textwidth}
      \begin{block}{Code}
        \texttt{car1 = Car("Honda", "Civic", 2020)}
      \end{block}

      \begin{itemize}
        \item The class call constructs a \texttt{Car} instance.
        \item Assignment binds the name \texttt{car1} to that resulting object.
        \item Access through \texttt{car1} follows that reference to the instance.
      \end{itemize}
    \end{column}
    \begin{column}{0.43\textwidth}
      % Intended visual: variable name car1 pointing to a Car instance with three attributes.
      \lectureimage[2.22in]{name-reference-object.png}
    \end{column}
  \end{columns}
\end{frame}

\begin{frame}{Identity Asks Whether Two References Reach the Same Object}
  \begin{cpdefinition}
    \textbf{Object identity} distinguishes one particular object from every other object, even when their current state happens to look the same.
  \end{cpdefinition}

  \begin{block}{Example}
    \texttt{a = Car("Honda", "Civic", 2020)}\\
    \texttt{b = Car("Honda", "Civic", 2020)}\\
    \texttt{a is b} $\rightarrow$ \texttt{False}
  \end{block}

  \begin{alertblock}{Next lecture}
    We will use identity to reason about aliasing when two names actually refer to the same instance.
  \end{alertblock}
\end{frame}

\sectionframe{Part 6: Objects as a Data Representation}

\begin{frame}{Dictionary and Object Can Represent Similar Information}
  \begin{columns}[T,onlytextwidth]
    \begin{column}{0.47\textwidth}
      \begin{block}{Dictionary}
        \texttt{car = \{}\\
        \hspace*{0.15in}\texttt{"make": "Honda",}\\
        \hspace*{0.15in}\texttt{"model": "Civic",}\\
        \hspace*{0.15in}\texttt{"year": 2020}\\
        \texttt{\}}\\[0.08in]
        \texttt{car["year"]}
      \end{block}
    \end{column}
    \begin{column}{0.47\textwidth}
      \begin{block}{Object}
        \texttt{car = Car(}\\
        \hspace*{0.15in}\texttt{"Honda",}\\
        \hspace*{0.15in}\texttt{"Civic",}\\
        \hspace*{0.15in}\texttt{2020}\\
        \texttt{)}\\[0.08in]
        \texttt{car.year}
      \end{block}
    \end{column}
  \end{columns}
\end{frame}

\begin{frame}{Object vs. Dictionary: What Changes?}
  \begin{center}
    \renewcommand{\arraystretch}{1.30}
    \begin{tabular}{l|l|l}
      \textbf{Question} & \textbf{Dictionary} & \textbf{Class instance} \\\hline
      field name & key & attribute \\
      access & \texttt{data["year"]} & \texttt{car.year} \\
      structure & convention among keys & described by class \\
      create another & new dictionary & call class again \\
      future behavior & separate functions & methods can belong to class
    \end{tabular}
  \end{center}

  \begin{keyidea}
    Dictionaries remain useful. Classes become useful when the program benefits from defining a reusable kind of entity with a known structure and, eventually, associated behavior.
  \end{keyidea}
\end{frame}

\sectionframe{Part 7: Connecting to Python You Already Know}

\begin{frame}{You Have Already Used Objects}
  Python's built-in data types also produce objects.

  \begin{columns}[T,onlytextwidth]
    \begin{column}{0.47\textwidth}
      \begin{block}{List}
        \texttt{values = [10, 20, 30]}\\
        \texttt{values.append(40)}
      \end{block}
    \end{column}
    \begin{column}{0.47\textwidth}
      \begin{block}{String}
        \texttt{text = "hello"}\\
        \texttt{text.upper()}
      \end{block}
    \end{column}
  \end{columns}

  \begin{itemize}
    \item We previously concentrated on the sequence behavior we needed.
    \item Now we have enough vocabulary to say that lists and strings are objects with methods provided by their classes.
  \end{itemize}
\end{frame}

\begin{frame}{Dot Notation Has Two Closely Related Roles}
  \begin{columns}[T,onlytextwidth]
    \begin{column}{0.47\textwidth}
      \begin{block}{Attribute access}
        \texttt{car1.year}
      \end{block}
      Retrieve data associated with an object.
    \end{column}
    \begin{column}{0.47\textwidth}
      \begin{block}{Method call}
        \texttt{values.append(40)}
      \end{block}
      Invoke behavior associated with an object.
    \end{column}
  \end{columns}

  \begin{keyidea}
    Today we defined our own attributes. Next lecture we define our own methods and connect method calls to particular instances.
  \end{keyidea}
\end{frame}

\begin{frame}{Design Prompt: What Should a \texttt{Student} Store?}
  \begin{activity}
    We want one value to represent one student in a course. What attributes belong to that object?
  \end{activity}

  Possible design:

  \begin{block}{Class}
    \texttt{class Student:}\\
    \hspace*{0.20in}\texttt{def \_\_init\_\_(self, name: str, major: str, gpa: float):}\\
    \hspace*{0.40in}\texttt{self.name = name}\\
    \hspace*{0.40in}\texttt{self.major = major}\\
    \hspace*{0.40in}\texttt{self.gpa = gpa}
  \end{block}
\end{frame}

\begin{frame}{The Data Model Should Match the Concept}
  A good class does not simply collect every value a program might ever use.

  \begin{itemize}
    \item Attributes should describe the state of the conceptual entity.
    \item Temporary calculations usually do not need to become attributes.
    \item Two attributes should not redundantly store the same information unless the design requires it.
    \item Attribute names should communicate what each component means.
  \end{itemize}

  \begin{exampleblockcp}
    A student's \texttt{name}, \texttt{major}, and \texttt{gpa} describe student state. A temporary variable such as \texttt{new\_gpa} used during a calculation does not automatically belong in the object.
  \end{exampleblockcp}
\end{frame}

\begin{frame}{Common Misconception: Attributes Are Local Variables}
  Inside \texttt{\_\_init\_\_}:

  \begin{block}{Different roles}
    \texttt{make} --- parameter/local name for this call\\
    \texttt{self.make} --- attribute stored on the instance
  \end{block}

  After \texttt{\_\_init\_\_} returns, the parameter name \texttt{make} is gone with that call frame, but \texttt{self.make} remains part of the object's state.

  \begin{keyidea}
    Parameters help initialize the object; attributes preserve the resulting state beyond the initialization call.
  \end{keyidea}
\end{frame}

\begin{frame}{Trace Checklist for Object Construction}
  When you encounter:

  \begin{block}{Statement}
    \texttt{name = ClassName(argument1, argument2, ...)}
  \end{block}

  trace it in this order:

  \begin{enumerate}
    \item evaluate the argument expressions;
    \item construct a new instance;
    \item bind initialization parameters;
    \item let \texttt{self} refer to the new instance;
    \item execute the attribute assignments;
    \item finish the class-call expression with the initialized instance;
    \item bind the result to the variable name.
  \end{enumerate}
\end{frame}

\begin{frame}{Where We Stop Today}
  By the end of Lecture 28, you should be able to:

  \begin{itemize}
    \item explain why an object is useful for compound data;
    \item distinguish a class from an instance;
    \item read a simple class containing \texttt{\_\_init\_\_};
    \item trace construction of a new instance;
    \item identify parameter names versus instance attributes;
    \item access and update attributes with dot notation;
    \item draw names as references to distinct object instances;
    \item compare an object representation with a dictionary representation.
  \end{itemize}
\end{frame}

\begin{frame}{Takeaways}
  \begin{enumerate}
    \item Compound data represents one conceptual value using several related components.
    \item A class defines a reusable kind of object; an instance is one particular object of that class.
    \item Calling a class constructs a new instance and \texttt{\_\_init\_\_} establishes its initial state.
    \item Attributes are named components of object state and are accessed with dot notation.
    \item Separate constructor calls produce separate object identities, even when their state looks the same.
    \item The reference model developed for lists generalizes naturally to user-defined objects.
  \end{enumerate}

  \begin{block}{Next: Objects II}
    Methods connect behavior to specific instances; aliasing and collections determine how those instances interact with the rest of a program.
  \end{block}
\end{frame}

\end{document}
